\chapter{Khung portal --- Đăng nhập, Trang chủ, Cửa hàng}

Chương này gồm \textbf{25 đường dẫn} gom thành \textbf{13 màn}, nằm trong hai mô-đun:

\begin{itemize}[leftmargin=*,itemsep=1pt]
  \item \rt{wujia_portal_layout} (11 đường dẫn) --- \emph{cái khung}: đăng nhập, quên và đổi
        mật khẩu, hồ sơ cá nhân, đổi ngôn ngữ. Phần này cố ý \textbf{không biết gì} về cửa
        hàng, đơn hàng hay công nợ, để màn đăng nhập vẫn chạy được kể cả khi các phân hệ
        nghiệp vụ chưa bật.
  \item \rt{wujia_portal_base} (14 đường dẫn) --- \emph{cái nền ghép}: Trang chủ tổng hợp số
        liệu từ các phân hệ, việc chọn cửa hàng đang làm việc, và các màn thông tin cửa hàng.
\end{itemize}

\textbf{Năm nguyên tắc chi phối toàn bộ portal}, mô tả một lần ở đây, các chương sau chỉ trỏ về:

\begin{enumerate}[leftmargin=*]
  \item \textbf{Cửa hàng đang làm việc nằm trong cookie, không nằm trong tài khoản.} Trình duyệt
        giữ \rt{wujia_active_franchise_id} trong \textbf{30 ngày}. Mỗi lần mở trang, hệ thống
        kiểm lại: cửa hàng trong cookie có nằm trong danh sách cửa hàng tài khoản còn hiệu lực
        không --- không thì coi như chưa chọn. Ai chỉ thuộc \textbf{một} cửa hàng thì hệ thống
        tự chọn, không hỏi. Hệ quả: đổi máy hoặc xoá dữ liệu trình duyệt là phải chọn lại.
  \item \textbf{Chưa chọn cửa hàng thì dữ liệu là TỔNG của mọi cửa hàng.} Đây là điểm dễ hiểu
        nhầm nhất. Với người thuộc nhiều cửa hàng mà chưa bấm chọn, hệ thống \emph{không} chặn,
        mà lấy dữ liệu của \textbf{tất cả} cửa hàng họ thuộc. Trang chủ vì thế có thể hiện con
        số cộng gộp của 3 cửa hàng.
  \item \textbf{Vai trò có ba bậc: Nhân viên $<$ Quản lý $<$ Chủ tiệm.} Bậc dùng để so sánh
        ``ít nhất phải là\ldots''. Một tài khoản thuộc nhiều cửa hàng thì vai trò được lấy là
        \textbf{bậc cao nhất} trong các cửa hàng đang xét. Rất nhiều màn hiện nay \emph{không}
        phân biệt vai trò --- chỗ nào có phân biệt sẽ được nói rõ ở mục 2 của màn đó.
  \item \textbf{Hầu hết dữ liệu được đọc bằng quyền hệ thống.} Nghĩa là phân quyền chuẩn của
        Odoo bị bỏ qua, và việc ``ai thấy gì'' \textbf{hoàn toàn} do điều kiện lọc trong từng
        màn quyết định. Nếu một điều kiện lọc viết thiếu, dữ liệu cửa hàng khác sẽ lọt ra --- đó
        là lý do tài liệu này liệt kê từng điều kiện một.
  \item \textbf{Giờ hiển thị đổi theo múi giờ của TÀI KHOẢN.} Hệ thống lưu giờ chuẩn quốc tế,
        khi hiện lên mới quy đổi theo ô \emph{Múi giờ} trong tài khoản người dùng; ô trống thì
        dùng \emph{Asia/Ho\_Chi\_Minh}. Không phải theo cửa hàng, cũng không phải theo khu vực.
\end{enumerate}

\section{Màn Đăng nhập}
\rt{/portal/login} · \rt{/portal/logout} · \rt{/portal/redirect}
\medskip

\mvao

Là màn đầu tiên. Mọi đường dẫn portal khi chưa đăng nhập đều bị đẩy về đây.
\rt{/portal/logout} là nút Đăng xuất; \rt{/portal/redirect} \textbf{không phải màn}, mà là
trạm trung chuyển ngay sau khi đăng nhập thành công.

\mai

\begin{itemize}[leftmargin=*]
  \item Màn này \textbf{không yêu cầu đăng nhập} (đương nhiên). Đã đăng nhập rồi mà mở lại
        \rt{/portal/login} thì bị đẩy thẳng qua trạm trung chuyển.
  \item Trạm trung chuyển phân luồng theo \textbf{loại tài khoản}: tài khoản \emph{nội bộ}
        (nhân sự Ngô Gia) đi vào giao diện quản trị \rt{/odoo}; tài khoản \emph{cửa hàng} đi
        vào \rt{/portal}. Trên UAT: \uat{1} tài khoản nội bộ, \uat{6} tài khoản cửa hàng.
  \item Nếu địa chỉ gọi tới có kèm tham số ``sau khi vào thì đi đâu'', hệ thống đi theo tham số
        đó \textbf{trước}, không qua trạm phân luồng.
\end{itemize}

\mhien

\begin{hienbang}
Ô tài khoản, ô mật khẩu &
  \rt{res.users} &
  Không truy vấn gì lúc hiện màn. Ô tài khoản được điền sẵn nếu phiên trước có ghi nhận. &
  --- & \uat{7} tài khoản đang có \\ \addlinespace
Nút đổi ngôn ngữ &
  \rt{res.lang} &
  Chỉ ngôn ngữ \textbf{đang bật} (\rt{active = true}) &
  --- & \uat{3} ngôn ngữ: Tiếng Việt, English, Thai \\
\end{hienbang}

\mloc

Hai ô: \emph{tài khoản} (chính là email đăng nhập) và \emph{mật khẩu}. Người chưa đăng nhập
vẫn đổi được ngôn ngữ ngay tại màn này --- lựa chọn đó được nhớ tạm và \textbf{ghi vào tài
khoản} ngay sau khi đăng nhập thành công.

\mkiem{Đăng nhập}

\begin{kiembang}
1 & Xác định đúng cơ sở dữ liệu cần đăng nhập &
    Máy chủ chỉ có một cơ sở dữ liệu nên bước này luôn qua \\
2 & Tài khoản và mật khẩu khớp &
    ``Wrong login/password''. Hệ thống \textbf{không} nói rõ sai tài khoản hay sai mật khẩu \\
3 & Tài khoản còn hoạt động &
    Thông báo do Odoo sinh ra (ví dụ tài khoản đã bị tắt) \\
\end{kiembang}

\textbf{Đáng lưu ý:} màn đăng nhập \textbf{không} có cơ chế giới hạn số lần thử của riêng
portal --- chỉ màn \emph{Quên mật khẩu} mới có (10 lần/giờ/địa chỉ mạng). Xem mục 7.

\mghi

\begin{itemize}[leftmargin=*]
  \item Mở một \emph{phiên làm việc} cho tài khoản. Không tạo bản ghi nghiệp vụ nào.
  \item Nếu người dùng đã đổi ngôn ngữ ở màn đăng nhập, hệ thống \textbf{ghi ngôn ngữ đó vào
        tài khoản} (\rt{res.users.lang}) --- đây là lần duy nhất màn đăng nhập ghi dữ liệu.
  \item Đăng xuất: đóng phiên, quay về màn đăng nhập. \textbf{Không} xoá cookie cửa hàng đang
        chọn --- lần sau đăng nhập lại vẫn còn nguyên cửa hàng cũ, kể cả khi người đăng nhập
        là tài khoản khác trên cùng máy (hệ thống vẫn kiểm lại quyền nên không lộ dữ liệu, chỉ
        là cửa hàng hiển thị có thể không như mong đợi).
\end{itemize}

\mhoi

\begin{ask}
\begin{enumerate}[leftmargin=*]
  \item \textbf{Chưa chặn thử mật khẩu hàng loạt ở màn đăng nhập.} Màn \emph{Quên mật khẩu} đã
        có chặn 10 lần/giờ/địa chỉ mạng, nhưng màn đăng nhập thì chưa. Anh có muốn đặt luật
        tương tự (ví dụ khoá tạm 5 phút sau 10 lần sai) không, và thông báo hiển thị thế nào?
  \item \textbf{Đăng xuất không xoá cửa hàng đang chọn.} Máy dùng chung ở cửa hàng, người sau
        đăng nhập sẽ thấy tên cửa hàng của người trước cho tới khi bấm đổi. Có cần xoá khi
        đăng xuất không?
\end{enumerate}
\end{ask}

\section{Màn Quên mật khẩu}
\rt{/portal/forgot-pass}
\medskip

\mvao

Liên kết ``Quên mật khẩu?'' ngay dưới ô mật khẩu ở màn đăng nhập.

\mai

Không cần đăng nhập. Ai cũng mở được.

\mhien

Một ô nhập \emph{email}, một nút gửi. Không hiển thị dữ liệu nào từ hệ thống.

\mloc

Chỉ ô email.

\mkiem{Gửi yêu cầu}

\begin{kiembang}
1 & Địa chỉ mạng chưa gửi quá \textbf{10 lần trong 1 giờ} &
    ``Quá nhiều yêu cầu. Vui lòng thử lại sau.'' \\
2 & Có nhập email &
    ``Vui lòng nhập email.'' \\
3 & \emph{(không kiểm email có tồn tại hay không --- xem giải thích dưới)} &
    --- \\
\end{kiembang}

\textbf{Cố ý không báo email có tồn tại hay không.} Dù email đúng hay sai, người dùng luôn nhận
đúng một câu: \emph{``Nếu email tồn tại trong hệ thống, hướng dẫn đặt lại mật khẩu đã được gửi
tới hộp thư của bạn.''} Đây là chủ đích --- nếu báo ``email không tồn tại'' thì kẻ xấu có thể
dò ra danh sách email đang dùng hệ thống. Lỗi thật (email không có, hoặc máy chủ thư chưa cấu
hình) chỉ được ghi vào nhật ký hệ thống cho kỹ thuật xem.

\mghi

Sinh một \textbf{mã dùng một lần} gắn vào tài khoản và gửi email chứa liên kết đặt lại. Nếu
máy chủ gửi thư chưa được cấu hình thì \textbf{không có gì được gửi đi} nhưng người dùng vẫn
thấy câu thông báo thành công.

\mhoi

\begin{ask}
Vì màn luôn báo thành công, người dùng nhập nhầm email sẽ ngồi chờ một email không bao giờ tới,
và Ngô Gia cũng không biết. Anh muốn giữ nguyên (an toàn hơn), hay thêm dòng hướng dẫn kiểu
``không nhận được sau 5 phút, vui lòng liên hệ hotline\ldots''?
\end{ask}

\section{Màn Đặt lại mật khẩu}
\rt{/portal/reset_password}
\medskip

\mvao

Bấm liên kết trong email đặt lại mật khẩu. Liên kết có kèm \textbf{mã dùng một lần}.

\mai

Không cần đăng nhập, nhưng \textbf{bắt buộc có mã hợp lệ}. Vào thẳng đường dẫn mà không có mã
thì bị đẩy về màn đăng nhập. Mã hết hạn hoặc đã dùng cũng bị đẩy về.

\mhien

Tên tài khoản (lấy ra từ chính mã), ô \emph{mật khẩu mới}, ô \emph{nhập lại mật khẩu}.

\mloc

Hai ô mật khẩu.

\mkiem{Đặt lại mật khẩu}

\begin{kiembang}
1 & Liên kết có mã & Đẩy về màn đăng nhập \\
2 & Hai ô mật khẩu khớp nhau & ``Mật khẩu xác nhận không khớp.'' \\
3 & Mật khẩu \textbf{tối thiểu 8 ký tự} & ``Mật khẩu tối thiểu 8 ký tự.'' \\
\end{kiembang}

\textbf{Chỉ có luật độ dài 8 ký tự} --- không bắt buộc chữ hoa, chữ số hay ký tự đặc biệt.

\mghi

Ghi mật khẩu mới vào tài khoản, huỷ hiệu lực của mã, rồi đưa người dùng về màn đăng nhập kèm
thông báo thành công. \textbf{Không} tự đăng nhập hộ.

\mhoi

\begin{ask}
Luật mật khẩu hiện chỉ là ``từ 8 ký tự''. Anh có yêu cầu thêm (hoa/số/ký tự đặc biệt, không
trùng mật khẩu cũ, buộc đổi sau N tháng) không? Lưu ý luật này phải áp \textbf{đồng thời} cho
màn Đổi mật khẩu ở mục 4.7.
\end{ask}

\section{Màn Trang chủ}
\rt{/portal}
\medskip

\mvao

Vào thẳng sau khi đăng nhập, hoặc bấm biểu tượng nhà trên thanh điều hướng.

\mai

\begin{itemize}[leftmargin=*]
  \item Phải đăng nhập.
  \item \textbf{Phạm vi dữ liệu}: nếu đã chọn cửa hàng thì \textbf{chỉ} cửa hàng đó; chưa chọn
        thì \textbf{cộng gộp mọi cửa hàng} tài khoản thuộc. Trên UAT có \uat{2}/7 tài khoản
        thuộc nhiều hơn một cửa hàng (một tài khoản 3 cửa hàng, một tài khoản 2 cửa hàng), nên
        hiện tượng cộng gộp này BA test bằng hai tài khoản đó.
  \item Tài khoản \textbf{không thuộc cửa hàng nào}: mọi con số về 0, mọi danh sách rỗng ---
        màn vẫn mở được, không báo lỗi.
  \item \textbf{Hộp chọn cửa hàng} tự bật lên khi tài khoản thuộc \textbf{từ 2 cửa hàng trở
        lên} \emph{và} cookie đang trống.
  \item \textbf{Vai trò không giới hạn gì ở màn này} --- Nhân viên nhìn thấy đúng bằng Chủ tiệm,
        kể cả các con số về đơn hàng và đổi trả.
  \item Toàn bộ số liệu đọc bằng quyền hệ thống; giới hạn duy nhất là điều kiện
        \rt{franchise_id} trong bảng dưới.
\end{itemize}

\mhien

\textbf{Bốn thẻ số liệu} ở đầu màn:

\begin{hienbang}
Thẻ \emph{Thông báo chưa đọc} &
  \rt{wujia.notification} trừ đi \rt{wujia.notification.read} &
  Số thông báo đã công bố (\rt{is_published_portal = true}) gửi cho \textbf{mọi cửa hàng}
  (\rt{franchise_ids} trống) \textbf{hoặc} đúng cửa hàng đang xét; rồi \textbf{trừ} đi số lượt
  người này đã đánh dấu đọc. Không âm. &
  --- &
  \uat{2} thông báo đã công bố, cả \uat{2} đều gửi toàn hệ \\ \addlinespace
Thẻ \emph{Đơn hàng 30 ngày} &
  \rt{sale.order} &
  Đơn của cửa hàng đang xét, ngày đặt trong \textbf{30 ngày} gần nhất, trạng thái khác
  \emph{Đã huỷ}. &
  --- &
  \uat{31} đơn chưa huỷ trên toàn hệ \\ \addlinespace
Thẻ \emph{Đơn chờ xác nhận} &
  \rt{sale.order} &
  Đơn của cửa hàng đang xét, trạng thái \emph{Nháp} hoặc \emph{Đã gửi}
  (\rt{state in ('draft','sent')}). Không giới hạn thời gian. &
  --- &
  \uat{17} đơn \\ \addlinespace
Thẻ \emph{Yêu cầu đổi trả đang mở} &
  \rt{wujia.return.request} &
  Yêu cầu của cửa hàng đang xét, trạng thái \emph{Đã gửi}, \emph{Đang xử lý} hoặc
  \emph{Đã duyệt} (\rt{state in ('submitted','processing','approved')}). &
  --- &
  \uat{7} yêu cầu \\
\end{hienbang}

\textbf{Các khối danh sách} bên dưới --- mọi khối đều cắt ở \textbf{2 dòng}, xem đầy đủ phải
bấm ``Xem tất cả'':

\begin{hienbang}
Thông báo mới nhất &
  \rt{wujia.notification} &
  Đã công bố \textbf{và} (gửi toàn hệ \textbf{hoặc} đúng cửa hàng đang xét) &
  Ngày công bố mới nhất trước; 2 dòng &
  \uat{2} thông báo \\ \addlinespace
Đơn hàng gần đây &
  \rt{sale.order} &
  Đơn của cửa hàng đang xét, trạng thái khác \emph{Đã huỷ} &
  Ngày đặt mới nhất trước; 2 dòng &
  \uat{31} đơn \\ \addlinespace
Yêu cầu đổi trả gần đây &
  \rt{wujia.return.request} &
  Của cửa hàng đang xét, \textbf{loại bỏ} \emph{Từ chối} và \emph{Đã huỷ}
  (\rt{state not in ('rejected','cancelled')}) --- lưu ý bộ trạng thái này
  \textbf{khác} bộ của thẻ số ở trên. &
  Ngày gửi mới nhất trước; 2 dòng &
  \uat{12} yêu cầu lọt bộ lọc này \\ \addlinespace
Sản phẩm đặt nhiều nhất &
  \rt{sale.order.line} &
  Dòng hàng thuộc đơn của cửa hàng đang xét, đơn đã \emph{Xác nhận} hoặc \emph{Hoàn tất}
  (\rt{state in ('sale','done')}), ngày đặt trong \textbf{90 ngày}. Gộp theo sản phẩm, cộng số
  lượng và thành tiền. &
  Số lượng nhiều nhất trước; 5 dòng &
  --- \\ \addlinespace
Giao hàng sắp tới \note{(bản điện thoại)} &
  \rt{stock.picking.batch} &
  Chuyến có phiếu giao thuộc cửa hàng đang xét, trạng thái \textbf{chưa hoàn thành}, và (khởi
  hành từ đầu hôm nay trở đi \textbf{hoặc} chưa đặt lịch \textbf{hoặc} đang bốc/đang chạy). &
  Giờ khởi hành sớm nhất trước; 2 dòng &
  \uat{1} chuyến chưa hoàn thành \\ \addlinespace
Bài kiến thức \note{(bản điện thoại)} &
  \rt{wujia.knowledge.article} &
  Đã công bố lên portal. \textbf{Không} lọc theo cửa hàng --- mọi cửa hàng thấy như nhau. &
  Ngày công bố mới nhất trước; 2 dòng &
  \uat{21} bài \\ \addlinespace
Dải chào ở đầu màn \note{(bản điện thoại)} &
  \rt{wujia.franchise.management} + \rt{wujia.franchise.member} &
  Tên cửa hàng đang chọn và vai trò của chính người đang xem tại cửa hàng đó. Chưa chọn cửa
  hàng thì để trống. &
  --- &
  \uat{4} cửa hàng, \uat{10} thành viên còn hiệu lực \\ \addlinespace
Dải khung giờ đặt hàng \note{(bản điện thoại)} &
  \rt{wujia.order.window} + tham số chung &
  Theo \textbf{khu vực của cửa hàng đang chọn}. Ba trạng thái: \emph{Đặt hàng 24/7} (khung giờ
  tắt), \emph{Đang mở} kèm thời gian còn lại, \emph{Đã đóng} kèm giờ mở lại. Luật đầy đủ ở
  \textbf{mục 3.7}. &
  --- &
  \uat{1} khung riêng / 2 khu vực \\
\end{hienbang}

\textbf{Ký hiệu tiền trên Trang chủ.} Số tiền gộp từ các đơn (bảng \emph{Sản phẩm đặt nhiều
nhất}, tổng tiền chuyến giao) mang ký hiệu tiền tệ \textbf{của chính các đơn đó} --- lấy từ bảng
giá gắn với khách hàng của cửa hàng, không phải tiền tệ của công ty. Nếu các đơn được gộp lại
đang dùng \textbf{nhiều loại tiền khác nhau} thì con số cộng gộp vốn đã vô nghĩa, hệ thống rơi
về tiền tệ công ty.

\mloc

Trang chủ \textbf{không có bộ lọc}. Thứ duy nhất người dùng đổi được là \emph{cửa hàng đang làm
việc} (mục 4.5) --- đổi cửa hàng thì toàn bộ số liệu trên màn đổi theo.

\mkiem{một thẻ số hoặc ``Xem tất cả''}

Không kiểm gì --- đây chỉ là liên kết sang màn danh sách của phân hệ tương ứng.

\mghi

\textbf{Trang chủ không ghi gì vào hệ thống.} Mở màn bao nhiêu lần cũng không tạo, không sửa
bản ghi nào; kể cả khối Thông báo cũng không tự đánh dấu là đã đọc.

\mhoi

\begin{ask}
\begin{enumerate}[leftmargin=*]
  \item \textbf{Số ``Thông báo chưa đọc'' ở Trang chủ có thể lệch với số trên màn Thông báo.}
        Ba khác biệt trong cách đếm: (a) Trang chủ \textbf{không} loại thông báo \emph{đã hết
        hiệu lực}, màn Thông báo thì có; (b) Trang chủ \textbf{không} loại thông báo có ngày
        công bố ở tương lai; (c) Trang chủ trừ \textbf{mọi} lượt đánh dấu đọc của người đó,
        trong khi trạng thái đọc được lưu \textbf{theo từng cửa hàng} --- người làm ở 2 cửa hàng
        đọc cùng một thông báo ở cả hai nơi sẽ bị trừ 2. Trên UAT tài khoản
        \emph{Administrator} thuộc 2 cửa hàng và đang có \uat{3} lượt đọc trên \uat{2} thông
        báo, nên con số bị kẹp về 0. Anh chốt giúp: lấy cách đếm của màn Thông báo làm chuẩn
        cho cả Trang chủ, đúng không?
  \item \textbf{Thẻ ``Đổi trả đang mở'' bỏ sót trạng thái \emph{Đang xét}.} Trạng thái
        \rt{reviewing} có thật trong hệ thống nhưng không nằm trong bộ đếm
        (\rt{submitted}, \rt{processing}, \rt{approved}). Trên UAT hiện chưa có yêu cầu nào ở
        trạng thái này nên chưa ai thấy. Có phải \emph{Đang xét} cũng là ``đang mở'' không?
  \item \textbf{Thẻ số và danh sách bên dưới dùng hai bộ trạng thái khác nhau.} Thẻ đếm
        \emph{Đã gửi + Đang xử lý + Đã duyệt}; danh sách hiện \emph{mọi trạng thái trừ Từ chối
        và Đã huỷ} --- nghĩa là danh sách \textbf{có} cả phiếu \emph{Nháp} và \emph{Hoàn thành}
        mà thẻ không đếm. Trên UAT: thẻ ra \uat{7}, danh sách lọc ra \uat{12}. Anh muốn thống
        nhất theo bộ nào?
  \item \textbf{Mọi khối danh sách chỉ hiện 2 dòng.} Con số 2 đang cố định trong code, không
        cấu hình được. Trên màn hình máy tính rộng, 2 dòng trông khá trống. Anh muốn giữ 2, hay
        đổi (ví dụ 5 dòng cho máy tính, 2 dòng cho điện thoại)?
  \item \textbf{Chưa chọn cửa hàng thì các con số là tổng của mọi cửa hàng.} Người quản lý nhiều
        cửa hàng mở Trang chủ lần đầu sẽ thấy số cộng gộp mà không có nhãn nào nói điều đó. Anh
        muốn (a) giữ nguyên nhưng thêm dòng ``Đang xem: tất cả 3 cửa hàng'', hay (b) bắt buộc
        chọn cửa hàng mới cho xem số?
  \item \textbf{Khối ``Bài kiến thức'' không lọc theo cửa hàng.} Mọi cửa hàng thấy chung
        \uat{21} bài. Nếu sau này có tài liệu riêng theo khu vực/cấp cửa hàng thì cần bổ sung
        điều kiện lọc --- anh xác nhận trước mắt là dùng chung?
\end{enumerate}
\end{ask}

\section{Chọn cửa hàng đang làm việc}
\rt{/portal/franchise/switch} · \rt{/portal/franchise/clear}
\medskip

\mvao

Bấm tên cửa hàng trên thanh trên cùng, hoặc chọn trong hộp tự bật ở Trang chủ. Hai đường dẫn
này \textbf{không phải màn} --- chúng chỉ nhận lệnh rồi đưa người dùng quay lại trang cũ.

\mai

Phải đăng nhập. Chỉ chọn được cửa hàng mà tài khoản đang có thành viên \textbf{còn hiệu lực}.

\mhien

Danh sách cửa hàng trong hộp chọn lấy từ \rt{wujia.franchise.management}, giới hạn đúng những
cửa hàng tài khoản thuộc về.

\mloc

Người dùng chọn đúng một cửa hàng, hoặc bấm ``Bỏ chọn''.

\mkiem{Chọn cửa hàng}

\begin{kiembang}
1 & Mã cửa hàng gửi lên là số & Đưa về Trang chủ, không đổi gì \\
2 & Mã đó nằm trong danh sách cửa hàng tài khoản còn hiệu lực &
    Đưa về Trang chủ, \textbf{không có thông báo lỗi} --- người dùng chỉ thấy cửa hàng không
    đổi \\
3 & Địa chỉ quay lại phải là \textbf{đường dẫn nội bộ} &
    Địa chỉ trỏ ra ngoài (trang web khác) bị thay bằng \rt{/portal}; đây là lớp chặn chuyển
    hướng sang trang giả mạo \\
\end{kiembang}

\mghi

\begin{itemize}[leftmargin=*]
  \item Ghi \textbf{cookie} \rt{wujia_active_franchise_id} trên máy người dùng, hạn \textbf{30
        ngày}. \emph{Không} ghi gì vào cơ sở dữ liệu --- đây là lý do đổi máy là mất lựa chọn.
  \item ``Bỏ chọn'' xoá cookie và đưa về Trang chủ.
  \item Cookie \textbf{đọc được bằng mã trình duyệt} (có chủ đích, để hiện tên cửa hàng ở thanh
        trên mà không phải hỏi máy chủ).
\end{itemize}

\mhoi

\begin{ask}
\begin{enumerate}[leftmargin=*]
  \item \textbf{Chọn cửa hàng không được phép thì im lặng.} Người dùng bấm vào cửa hàng vừa bị
        gỡ quyền sẽ chỉ thấy trang tải lại mà không đổi gì, không có câu nào giải thích. Anh
        muốn thêm thông báo ``Bạn không còn quyền tại cửa hàng này'' không?
  \item \textbf{Cửa hàng đang chọn lưu trên trình duyệt.} Dùng máy khác, trình duyệt ẩn danh
        hoặc xoá dữ liệu duyệt web là phải chọn lại từ đầu. Anh xác nhận chấp nhận cách này,
        hay muốn nhớ theo tài khoản (lưu vào hồ sơ người dùng)?
\end{enumerate}
\end{ask}

\section{Màn Hồ sơ cá nhân}
\rt{/portal/profile} · \rt{/portal/profile/update} · \rt{/portal/profile/avatar/<int:user_id>}
\medskip

\mvao

Bấm ảnh đại diện ở góc trên bên phải $\rightarrow$ \emph{Hồ sơ}.
\rt{/portal/profile/update} là hành động Lưu; \rt{/portal/profile/avatar/<id>} là đường dẫn
\textbf{trả về ảnh đại diện} --- mọi chỗ trong portal hiện ảnh người dùng đều gọi qua đây.

\mai

\begin{itemize}[leftmargin=*]
  \item Phải đăng nhập. Mỗi người chỉ xem và sửa \textbf{hồ sơ của chính mình} --- không có
        đường dẫn nào cho xem hồ sơ người khác.
  \item \textbf{Ảnh đại diện thì khác}: xem được ảnh của \textbf{chính mình}, và ảnh của bất kỳ
        ai \textbf{cùng thuộc một cửa hàng} với mình. Người ngoài cửa hàng bị từ chối. Với dữ
        liệu UAT, tài khoản thuộc cửa hàng Cầu Giấy thấy được ảnh của \uat{5} người cùng cửa hàng
        đó (kể cả chính mình).
  \item Khối ``Cửa hàng / Vai trò'' trên màn chỉ hiện khi \textbf{đã chọn được cửa hàng}; người
        nhiều cửa hàng chưa chọn thì hai dòng này hiện dấu ``---''.
\end{itemize}

\mhien

\begin{hienbang}
Họ tên, email liên hệ, điện thoại, địa chỉ &
  \rt{res.partner} (liên hệ gắn với tài khoản) &
  Chính liên hệ của người đang đăng nhập &
  --- & --- \\ \addlinespace
Ảnh đại diện &
  \rt{res.users} &
  Bản ảnh cỡ 256 điểm ảnh. Trình duyệt được phép giữ lại \textbf{5 phút} nên đổi ảnh xong có
  thể phải tải lại trang mới thấy. &
  --- & --- \\ \addlinespace
Cửa hàng + Vai trò &
  \rt{wujia.franchise.member} &
  Thành viên \textbf{còn hiệu lực} của người này tại cửa hàng đang chọn. Chưa chọn và thuộc
  nhiều cửa hàng $\rightarrow$ để trống. &
  --- &
  \uat{10} thành viên còn hiệu lực (\uat{3} Chủ tiệm, \uat{3} Quản lý, \uat{4} Nhân viên) \\
\end{hienbang}

\mloc

Bốn ô nhập: \emph{họ tên}, \emph{email}, \emph{điện thoại}, \emph{địa chỉ}; và một ô chọn ảnh.

\mkiem{Lưu}

\begin{kiembang}
1 & Họ tên không được trống & ``Họ tên không được trống'' \\
2 & Điện thoại (nếu có nhập) đúng khuôn: 7--20 ký tự, chỉ số và các dấu \rt{+ ( ) . -} và
    khoảng trắng & ``Số điện thoại không hợp lệ'' \\
3 & \emph{(ảnh, nếu có chọn)} kích thước không quá \textbf{2MB} & ``Ảnh đại diện vượt quá 2MB'' \\
4 & \emph{(ảnh)} định dạng thuộc \textbf{PNG, JPEG, WEBP} & ``Định dạng ảnh không hỗ trợ'' \\
\end{kiembang}

\textbf{Email không bị kiểm khuôn} ở tầng này --- việc kiểm do Odoo làm khi ghi; sai thì hiện
thông báo do Odoo sinh ra.

\mghi

\begin{itemize}[leftmargin=*]
  \item Ghi \emph{họ tên}, \emph{điện thoại}, \emph{địa chỉ} vào liên hệ của tài khoản.
  \item \textbf{Email chỉ ghi vào liên hệ, KHÔNG đổi email đăng nhập.} Người dùng sửa email ở
        đây rồi lần sau vẫn phải đăng nhập bằng email cũ. Mỗi lần sửa đều được ghi vào nhật ký
        hệ thống.
  \item Ảnh (nếu có) ghi vào tài khoản ở cỡ gốc; Odoo tự sinh các cỡ nhỏ.
  \item Lưu xong hệ thống \textbf{chuyển hướng} về lại màn hồ sơ kèm câu ``Cập nhật thành công''
        --- cách này để bấm F5 không lưu lại lần nữa.
\end{itemize}

\mhoi

\begin{ask}
\begin{enumerate}[leftmargin=*]
  \item \textbf{Sửa email ở Hồ sơ không đổi được email đăng nhập} --- đây là điểm gần như chắc
        chắn gây hiểu nhầm. Anh muốn (a) đổi nhãn ô thành ``Email liên hệ'' và ghi chú rõ, hay
        (b) cho đổi luôn email đăng nhập (phải có bước xác nhận qua email mới)?
  \item \textbf{Ai cùng cửa hàng đều xem được ảnh đại diện của nhau.} Xác nhận đúng ý anh chứ?
  \item \textbf{Hồ sơ chưa cho sửa gì khác} --- không có ngày sinh, không có chức danh, không
        có liên kết tới hồ sơ nhân sự. Anh cần thêm trường nào không?
\end{enumerate}
\end{ask}

\section{Màn Đổi mật khẩu}
\rt{/portal/change-password}
\medskip

\mvao

Bấm ảnh đại diện góc trên bên phải $\rightarrow$ \emph{Đổi mật khẩu}.

\mai

Phải đăng nhập. Chỉ đổi được mật khẩu của chính mình. Không có đường dẫn nào cho Chủ tiệm đặt
lại mật khẩu cho nhân viên --- việc đó phải làm trong giao diện quản trị.

\mhien

Ba ô mật khẩu, kèm khối ``Cửa hàng / Vai trò'' giống màn Hồ sơ.

\mloc

\emph{Mật khẩu hiện tại}, \emph{mật khẩu mới}, \emph{nhập lại mật khẩu mới}.

\mkiem{Đổi mật khẩu}

\begin{kiembang}
1 & Có nhập cả mật khẩu cũ và mật khẩu mới & ``Vui lòng nhập đầy đủ.'' \\
2 & Hai ô mật khẩu mới khớp nhau & ``Mật khẩu xác nhận không khớp với mật khẩu mới.'' \\
3 & Mật khẩu mới \textbf{tối thiểu 8 ký tự} & ``Mật khẩu mới tối thiểu 8 ký tự.'' \\
4 & Mật khẩu mới \textbf{khác} mật khẩu cũ & ``Mật khẩu mới phải khác mật khẩu cũ.'' \\
5 & Mật khẩu hiện tại \textbf{đúng} &
    ``Mật khẩu hiện tại không đúng. Vui lòng kiểm tra lại.'' \\
\end{kiembang}

Lưu ý thứ tự: hệ thống chỉ kiểm mật khẩu cũ ở \textbf{bước cuối}, sau khi mật khẩu mới đã qua
hết các luật khuôn dạng.

\mghi

Ghi mật khẩu mới vào tài khoản. \textbf{Phiên đang mở vẫn giữ nguyên} --- người dùng không bị
đăng xuất, và các thiết bị khác đang đăng nhập cũng không bị đẩy ra.

\mhoi

\begin{ask}
\begin{enumerate}[leftmargin=*]
  \item \textbf{Đổi mật khẩu không đăng xuất các thiết bị khác.} Nếu ai đó đổi mật khẩu vì nghi
        bị lộ, thiết bị của kẻ kia \textbf{vẫn đang đăng nhập}. Anh có muốn đóng mọi phiên khác
        khi đổi mật khẩu không?
  \item Luật mật khẩu ở đây (8 ký tự, khác mật khẩu cũ) \textbf{không giống hoàn toàn} màn Đặt
        lại mật khẩu qua email (chỉ 8 ký tự, không kiểm khác cũ). Cần thống nhất --- anh chốt
        bộ luật chung.
\end{enumerate}
\end{ask}

\section{Đổi ngôn ngữ}
\rt{/portal/set-lang/<string:lang_code>}
\medskip

\mvao

Nút cờ/ngôn ngữ trên thanh trên cùng và ở màn đăng nhập. Không phải màn --- bấm xong quay lại
đúng trang đang đứng.

\mai

\textbf{Không cần đăng nhập} --- có chủ đích, để khách đổi được ngôn ngữ ngay ở màn đăng nhập.

\mhien

Danh sách ngôn ngữ lấy từ \rt{res.lang} với điều kiện \rt{active = true}. Trên UAT đang bật
\uat{3} ngôn ngữ: Tiếng Việt, English (US), Thai.

\mloc

Chọn một ngôn ngữ.

\mkiem{một ngôn ngữ}

\begin{kiembang}
1 & Mã ngôn ngữ có trong danh sách \textbf{đang bật} &
    Không đổi gì, vẫn quay lại trang cũ (im lặng bỏ qua) \\
2 & Địa chỉ quay lại là \textbf{đường dẫn nội bộ cùng máy chủ} &
    Bị thay bằng \rt{/portal} --- lớp chặn chuyển hướng sang trang giả mạo \\
\end{kiembang}

\mghi

\begin{itemize}[leftmargin=*]
  \item Cập nhật ngôn ngữ cho \textbf{phiên làm việc hiện tại} (có hiệu lực ngay ở trang kế).
  \item \textbf{Đã đăng nhập} $\rightarrow$ ghi luôn vào tài khoản (\rt{res.users.lang}), nên
        lần sau đăng nhập vẫn giữ ngôn ngữ đó.
  \item \textbf{Chưa đăng nhập} $\rightarrow$ chỉ nhớ tạm trong phiên, và được ghi vào tài
        khoản ngay sau khi đăng nhập thành công.
\end{itemize}

\mhoi

\begin{ask}
Portal đang bật \uat{3} ngôn ngữ, nghĩa là nút đổi ngôn ngữ hiện đủ 3 lựa chọn. Nếu bản dịch
tiếng Anh/tiếng Thái chưa đầy đủ thì người dùng chọn vào sẽ thấy màn lẫn lộn hai thứ tiếng. Anh
muốn tạm tắt bớt ngôn ngữ chưa dịch xong không?
\end{ask}

\section{Màn Danh sách cửa hàng của tôi}
\rt{/portal/franchises} · \rt{/my/franchises}
\medskip

\mvao

Menu tài khoản $\rightarrow$ \emph{Cửa hàng của tôi}. Hai đường dẫn cho \textbf{cùng một nội
dung}: bản \rt{/portal/\ldots} dùng giao diện Vuexy của dự án, bản \rt{/my/\ldots} dùng giao
diện portal mặc định của Odoo. Xem mục 7.

\mai

Phải đăng nhập. Chỉ hiện những cửa hàng mà tài khoản có thành viên \textbf{còn hiệu lực}.
Vai trò không giới hạn gì --- Nhân viên cũng thấy đủ.

\mhien

\begin{hienbang}
Danh sách cửa hàng &
  \rt{wujia.franchise.member} $\rightarrow$ \rt{wujia.franchise.management} &
  Thành viên của chính người đang đăng nhập, \rt{is_currently_valid = true}. ``Còn hiệu lực''
  nghĩa là: bản ghi chưa lưu trữ \textbf{và} đang làm việc \textbf{và} hôm nay nằm trong khoảng
  \emph{từ ngày} -- \emph{đến ngày}. &
  Theo thứ tự mặc định của bảng thành viên (cửa hàng, chủ chính trước, ngày bắt đầu mới trước)
  &
  \uat{10} bản ghi thành viên, \uat{0} bản hết hiệu lực \\ \addlinespace
Nhãn vai trò trên từng dòng &
  \rt{wujia.franchise.member.role} &
  \emph{Chủ tiệm} / \emph{Quản lý} / \emph{Nhân viên} &
  --- &
  \uat{3} / \uat{3} / \uat{4} \\
\end{hienbang}

\mloc

Không có bộ lọc, không phân trang. Một người thuộc bao nhiêu cửa hàng thì hiện bấy nhiêu dòng
(nhiều nhất trên UAT là \uat{3}).

\mkiem{một cửa hàng}

Không kiểm gì --- sang màn chi tiết.

\mghi

Không ghi gì.

\mhoi

\begin{ask}
\textbf{Đang có hai bộ màn song song cho cùng một nội dung}: \rt{/portal/franchises} (giao diện
Vuexy) và \rt{/my/franchises} (giao diện Odoo mặc định), tương tự ở màn chi tiết. Hai bộ này lấy
cùng dữ liệu nhưng khác giao diện, và \textbf{không} có liên kết nào trong portal trỏ tới bộ
\rt{/my/\ldots}. Anh cho biết: giữ cả hai (bộ \rt{/my/} để nhân sự Ngô Gia đối chiếu), hay gỡ
hẳn bộ \rt{/my/}? \note{Gỡ thì cần giữ chuyển hướng để liên kết cũ không chết.}
\end{ask}

\section{Màn Chi tiết cửa hàng}
\rt{/portal/franchises/<int:franchise_id>} · \rt{/my/franchises/<int:franchise_id>} ·
\rt{/portal/franchises/<int:franchise_id>/profile} ·
\rt{/my/franchises/<int:franchise_id>/members}
\medskip

\mvao

Bấm một dòng ở màn Danh sách cửa hàng. Đường dẫn thứ ba (\rt{/profile}) là bản \emph{hồ sơ đầy
đủ} của cửa hàng; đường dẫn thứ tư \textbf{không phải màn} mà là nguồn dữ liệu danh sách thành
viên cho giao diện gọi ngầm.

\mai

\begin{itemize}[leftmargin=*]
  \item Phải đăng nhập \textbf{và} có thành viên còn hiệu lực tại \textbf{đúng} cửa hàng đó.
        Không đạt $\rightarrow$ bị đưa về danh sách cửa hàng, \textbf{không báo lỗi} (cố ý
        không cho biết cửa hàng đó có tồn tại hay không).
  \item \textbf{Đây là điểm quan trọng về quyền riêng tư}: bước kiểm ở trên chỉ hỏi ``bạn có
        thuộc cửa hàng này không''. Sau khi qua được, danh sách \textbf{tất cả thành viên} của
        cửa hàng được đọc bằng quyền hệ thống. Nghĩa là \textbf{Nhân viên cũng xem được đầy đủ
        danh sách đồng nghiệp}, gồm tên, vai trò, số điện thoại, ngày bắt đầu và ngày kết thúc.
        Xem mục 7.
  \item Vai trò \textbf{không} giới hạn gì trong màn này.
\end{itemize}

\mhien

\begin{hienbang}
Khối thông tin cửa hàng &
  \rt{wujia.franchise.management} &
  Đúng cửa hàng được mở &
  --- &
  \uat{4} cửa hàng, cả \uat{4} đều \emph{Đang hoạt động} \\ \addlinespace
Khối thành viên &
  \rt{wujia.franchise.member} &
  Mọi thành viên của cửa hàng đó, \rt{is_currently_valid = true}. \textbf{Không} lọc theo người
  đang xem. &
  Theo thứ tự mặc định của bảng &
  Cầu Giấy \uat{5} người · TP HCM Quận 1 \uat{3} · Đống Đa \uat{1} · TEST \uat{1} \\ \addlinespace
Trạng thái cửa hàng &
  \rt{wujia.franchise.management.status} &
  Nhãn tiếng Việt cố định trong portal: \emph{Nháp} / \emph{Đang hoạt động} / \emph{Khoá} /
  \emph{Đã đóng} / \emph{Hết hạn} &
  --- &
  \uat{4}/4 \emph{Đang hoạt động} \\ \addlinespace
Hồ sơ đầy đủ \note{(đường dẫn /profile)} &
  \rt{wujia.franchise.management} + \rt{wujia.franchise.member} &
  Thêm các thông tin hợp đồng, khu vực, ngày bắt đầu/kết thúc nhượng quyền --- vẫn chỉ của cửa
  hàng người dùng thuộc về &
  --- & --- \\
\end{hienbang}

\mloc

Không có bộ lọc, không phân trang trên các màn này (khác với màn \emph{Thông tin cửa hàng} ở
mục 4.11 --- màn đó \textbf{có} phân trang).

\mkiem{một liên kết}

Không kiểm gì.

\mghi

Không ghi gì.

\mhoi

\begin{ask}
\begin{enumerate}[leftmargin=*]
  \item \textbf{Nhân viên xem được đầy đủ danh sách đồng nghiệp}, kèm số điện thoại và ngày
        hiệu lực. Anh xác nhận đúng mong muốn, hay muốn giới hạn: chỉ \emph{Chủ tiệm} và
        \emph{Quản lý} thấy đủ, còn \emph{Nhân viên} chỉ thấy tên và vai trò?
  \item \textbf{Bốn đường dẫn cho ba màn gần giống nhau.} \rt{/portal/franchises/<id>},
        \rt{/my/franchises/<id>} và \rt{/portal/franchises/<id>/profile} hiển thị phần lớn cùng
        thông tin, cộng thêm màn \emph{Thông tin cửa hàng} ở mục 4.11 nữa là \textbf{bốn} chỗ
        xem thông tin cửa hàng. Anh chốt giúp giữ lại \textbf{một} màn chính là màn nào; các
        đường dẫn còn lại sẽ chuyển hướng về đó.
\end{enumerate}
\end{ask}

\section{Màn Thông tin cửa hàng}
\rt{/portal/franchise-information}
\medskip

\mvao

Menu \emph{Thông tin cửa hàng} trên thanh điều hướng. Màn này luôn nói về \textbf{cửa hàng đang
chọn}, không nhận mã cửa hàng trên đường dẫn.

\mai

\begin{itemize}[leftmargin=*]
  \item Phải đăng nhập \textbf{và} đã chọn được cửa hàng. Chưa chọn $\rightarrow$ bị đưa về
        Trang chủ.
  \item Phải có thành viên còn hiệu lực tại cửa hàng đó.
  \item \textbf{Cổng chặn cứng}: cửa hàng \emph{bị khoá portal} (\rt{portal_locked = true})
        hoặc \textbf{không} ở trạng thái \emph{Đang hoạt động} (\rt{status != 'active'})
        $\rightarrow$ hiện màn ``bị khoá'' thay cho nội dung. Trên UAT hiện \uat{0} cửa hàng bị
        khoá và \uat{0} cửa hàng ngoài trạng thái \emph{Đang hoạt động}, nên \textbf{nhánh này
        chưa từng được kiểm thử bằng dữ liệu thật}.
  \item Vai trò không giới hạn gì.
\end{itemize}

\mhien

\begin{hienbang}
Khối hồ sơ cửa hàng &
  \rt{wujia.franchise.management} &
  Cửa hàng đang chọn &
  --- & --- \\ \addlinespace
Khối thành viên &
  \rt{wujia.franchise.member} &
  Thành viên của cửa hàng đang chọn, \rt{is_currently_valid = true} &
  Theo \textbf{vai trò} rồi mã bản ghi; \textbf{10} người mỗi trang &
  Cầu Giấy \uat{5} người --- chưa cửa hàng nào chạm ngưỡng phân trang \\ \addlinespace
Khối thành viên của chính tôi &
  \rt{wujia.franchise.member} &
  Bản ghi thành viên của người đang xem tại cửa hàng này &
  --- & --- \\
\end{hienbang}

\mloc

\begin{itemize}[leftmargin=*]
  \item \textbf{Số người mỗi trang} --- chỉ nhận \textbf{10, 20, 50}; gửi số khác thì lặng lẽ
        dùng 10.
  \item \textbf{Trang} --- vượt quá trang cuối thì tự kéo về trang cuối.
\end{itemize}

\mkiem{---}

Màn \textbf{chỉ xem}, không có nút ghi nào. Việc đề nghị sửa thông tin cửa hàng nằm ở phân hệ
\emph{Yêu cầu thông tin} (chương 7).

\mghi

Không ghi gì.

\mhoi

\begin{ask}
\begin{enumerate}[leftmargin=*]
  \item \textbf{Sắp xếp thành viên theo vai trò đang là thứ tự chữ cái của mã trạng thái}
        (\rt{manager} $\rightarrow$ \rt{owner} $\rightarrow$ \rt{staff}), tức ra \emph{Quản lý
        $\rightarrow$ Chủ tiệm $\rightarrow$ Nhân viên} --- không phải thứ tự cấp bậc. Anh muốn
        đổi thành \emph{Chủ tiệm $\rightarrow$ Quản lý $\rightarrow$ Nhân viên} không?
  \item \textbf{Nhánh ``cửa hàng bị khoá'' chưa có dữ liệu để test.} Nếu anh muốn nghiệm thu
        nhánh này, cần Ngô Gia khoá thử một cửa hàng trên UAT --- em không tự đổi dữ liệu.
\end{enumerate}
\end{ask}

\section{Bốn đường dẫn cũ đã đổi tên}
\rt{/my/branches} · \rt{/my/branches/<int:branch_id>} · \rt{/portal/branches} ·
\rt{/portal/branches/<int:branch_id>}
\medskip

Trước đây cửa hàng nhượng quyền được gọi là ``branch''. Khi đổi sang ``franchise'', bốn đường
dẫn cũ được \textbf{giữ lại làm đường chuyển hướng vĩnh viễn} sang đường dẫn mới tương ứng, để
liên kết cũ (trong email, tin nhắn, dấu trang của người dùng) không bị chết.

Bốn đường dẫn này \textbf{không hiện màn nào}, không đọc và không ghi dữ liệu. BA không cần lên
task cho chúng; chỉ cần biết là chúng tồn tại khi nhìn danh sách đường dẫn.

\section{Trang xem trước bộ giao diện}
\rt{/portal/_pc-preview}
\medskip

Trang \textbf{dành riêng cho đội phát triển}: bày sẵn mọi thành phần giao diện (nút, thẻ, bảng,
phân trang, nhãn trạng thái) để đối chiếu với bản thiết kế. \textbf{Chỉ nhân sự nội bộ mở
được} --- tài khoản cửa hàng vào sẽ bị đưa về Trang chủ. Không có liên kết nào trong portal trỏ
tới trang này, và nó không đọc dữ liệu nghiệp vụ nào.

\section{Quy ước dùng chung toàn portal}
\label{sec:quy-uoc-chung}

Những quy tắc dưới đây được viết \textbf{một lần duy nhất} trong mã nguồn và mọi màn của mọi
phân hệ đều dùng chung. Các chương sau sẽ trỏ về mục này thay vì lặp lại.

\subsection*{Phân trang}

\begin{itemize}[leftmargin=*]
  \item Số bản ghi mỗi trang \textbf{chỉ nhận đúng các giá trị có trong ô chọn} của màn đó
        (thường là 10 / 20 / 50, riêng màn Đặt hàng là 24 / 48 / 96). Gửi số khác --- kể cả sửa
        tay trên thanh địa chỉ --- thì hệ thống lặng lẽ dùng giá trị mặc định, \textbf{không
        báo lỗi}. Trần cứng là 100 bản ghi/trang.
  \item Số trang vượt quá trang cuối thì \textbf{tự kéo về trang cuối}.
  \item \textbf{Đổi trang luôn giữ nguyên mọi bộ lọc}: từ khoá, khoảng ngày, trạng thái đang
        chọn đều được mang theo. Đây là lý do các màn không còn hiện tượng ``bấm trang 2 là mất
        bộ lọc''.
  \item Thanh phân trang hiện dạng \emph{1 \ldots 4 5 6 \ldots 20} --- luôn có trang đầu, trang
        cuối, và một trang liền kề mỗi bên.
\end{itemize}

\subsection*{Ngày giờ và múi giờ}

\begin{itemize}[leftmargin=*]
  \item Hệ thống lưu giờ chuẩn quốc tế; khi hiện lên mới quy đổi theo ô \emph{Múi giờ} của
        \textbf{tài khoản người dùng}. Ô trống thì dùng \emph{Asia/Ho\_Chi\_Minh}. Múi giờ ghi
        sai (giá trị lạ) cũng rơi về mặc định chứ \textbf{không} làm lỗi trang.
  \item Trên UAT: \uat{1} tài khoản để trống múi giờ, \uat{6} tài khoản đặt giá trị khác
        \emph{Asia/Ho\_Chi\_Minh} --- cụ thể \emph{Asia/Saigon} (4 tài khoản) và
        \emph{Asia/Bangkok} (2 tài khoản). Cả ba giá trị này hiện đều là UTC+7 nên chưa gây
        lệch giờ; nhưng nếu sau này có người dùng ở múi giờ khác thì \textbf{mọi màn lọc theo
        ngày sẽ lệch theo tài khoản đó}, không theo cửa hàng.
  \item \textbf{Lọc theo khoảng ngày} được quy đổi đúng biên: ``từ ngày X'' tính từ 00:00 giờ
        địa phương, ``đến ngày Y'' tính hết 23:59:59 giờ địa phương.
  \item Khoảng ngày ngược (\emph{Từ ngày} sau \emph{Đến ngày}) dùng chung một câu báo trên mọi
        màn: \emph{``Từ ngày không được lớn hơn Đến ngày''}. Ngày sai định dạng thì tuỳ màn:
        có màn coi như bỏ lọc, có màn báo riêng.
\end{itemize}

\subsection*{Định dạng tiền}

\begin{itemize}[leftmargin=*]
  \item Một chỗ duy nhất sinh ra chuỗi tiền cho cả portal. Tiền Việt hiện dạng
        \emph{48.000 đ} (không phần lẻ, nhóm nghìn bằng dấu chấm); tiền có phần lẻ hiện dạng
        \emph{10.990,99 \$} --- số chữ số lẻ và ký hiệu đều lấy theo \textbf{loại tiền của
        chứng từ}, không cố định trong giao diện.
  \item Với các con số \textbf{cộng gộp nhiều đơn}, ký hiệu tiền lấy theo loại tiền của chính
        các đơn đó; trộn nhiều loại tiền thì rơi về tiền tệ công ty (và con số gộp khi đó vốn
        đã không có ý nghĩa).
\end{itemize}

\subsection*{Nhãn trạng thái (màu badge)}

\begin{itemize}[leftmargin=*]
  \item Màu của nhãn trạng thái \textbf{do chữ trên nhãn quyết định}, không do từng màn tự
        chọn --- nên cùng một chữ (ví dụ \emph{Đã xác nhận}) ra cùng một màu ở mọi màn, trên cả
        máy tính lẫn điện thoại. Đây là kết quả xử lý lỗi \emph{UI-STATUSBADGE-001}.
  \item Bảy bậc màu: \emph{trung tính} (chưa bắt đầu / đã khép lại) · \emph{thông tin} (mới, đã
        xác nhận) · \emph{chờ} (đang chờ bên khác) · \emph{đang chạy} · \emph{thành công} ·
        \emph{hỏng/dừng} · \emph{cần bạn làm gì đó}.
  \item Nhãn lạ chưa nằm trong bảng thì ra màu trung tính, \textbf{không bao giờ mất nhãn}.
  \item \textbf{Bản điện thoại có một số nhãn khác bản máy tính} theo đúng thiết kế đã duyệt ---
        ví dụ ticket hỗ trợ ở trạng thái chờ khách trả lời: điện thoại ghi \emph{``Có phản
        hồi''}, máy tính ghi \emph{``Chờ phản hồi''}. Đây là khác biệt \textbf{có chủ đích}
        theo bản thiết kế, xem mục 7 chương 7.
\end{itemize}

\subsection*{Tệp đính kèm}

\begin{itemize}[leftmargin=*]
  \item Mọi màn cho tải tệp lên (đổi trả, hỗ trợ, yêu cầu thông tin) dùng chung một bộ luật:
        tối đa \textbf{10 tệp}, mỗi tệp tối đa \textbf{5MB}, và \textbf{định dạng được kiểm ở
        máy chủ} chứ không tin khai báo của trình duyệt. Màn ảnh chỉ nhận PNG/JPEG/WEBP; màn tài
        liệu nhận thêm PDF.
  \item Tên tệp được làm sạch trước khi lưu (chống tên tệp độc hại).
  \item \textbf{Xem lại tệp đã tải lên}: hệ thống kiểm người xem có thuộc cửa hàng của chứng từ
        gốc không; không thuộc thì từ chối.
\end{itemize}

\subsection*{Vai trò}

Thang bậc \emph{Nhân viên (1) $<$ Quản lý (2) $<$ Chủ tiệm (3)}. Chỗ nào yêu cầu ``ít nhất
Quản lý'' thì Chủ tiệm đương nhiên qua. Người thuộc nhiều cửa hàng thì lấy \textbf{bậc cao
nhất}. Không đủ bậc $\rightarrow$ hệ thống từ chối thẳng với mã lỗi ``không có quyền'', không
phải chuyển hướng êm.

\subsection*{Giới hạn số lần gọi}

Chỉ \textbf{một} chỗ trong portal đang đặt giới hạn: màn \emph{Quên mật khẩu} --- 10 lần/giờ cho
mỗi địa chỉ mạng. Bộ đếm nằm trong bộ nhớ của tiến trình máy chủ, nghĩa là khởi động lại máy chủ
là bộ đếm về 0.
